\chapter{The quotient ring and its order}

\section{Multiplication on equivalence classes}

\begin{theorem}[Multiplication on surreal representatives]
  \label{thm:surreal-mul}
  \lean{Surreal.mul_isSurreal,Surreal.mul,Surreal.mul_congr}
  \leanok
  If \(a,b\) are surreal representatives, their game product is surreal.  Moreover, replacing
  either representative by an equivalent one does not change the equivalence class of the
  product.  Thus multiplication is a well-defined operation on representatives modulo
  \(\GameEq\).
\end{theorem}
\begin{proof}
  \leanok
  \uses{def:surreal-representative,def:is-surreal,def:game-mul,def:game-order,thm:conway-a,thm:game-mul-congr}
  By Definition~\ref{def:surreal-representative}, write
  \[
    a=(a_0,h_a),\qquad b=(b_0,h_b),
  \]
  where \(a_0,b_0\) are short games and \(h_a,h_b\) witness that they satisfy
  Definition~\ref{def:is-surreal}.  Theorem~\ref{thm:conway-a} shows that the
  product \(a_0b_0\), formed as in Definition~\ref{def:game-mul}, is again
  surreal.  Hence \((a_0b_0,h_{ab})\) is a surreal representative, where
  \(h_{ab}\) denotes the resulting witness.

  For alternative representatives
  \(a'=(a_1,h_{a_1})\) and \(b'=(b_1,h_{b_1})\), assume the underlying
  games satisfy \(a_0\GameEq a_1\) and \(b_0\GameEq b_1\) in the sense of
  Definition~\ref{def:game-order}.  Theorem~\ref{thm:game-mul-congr} then gives
  \[
    a_0b_0\GameEq a_1b_1.
  \]
  Thus the equivalence class of the product is independent of both chosen
  representatives.
\end{proof}
\begin{definition}[Multiplication on short surreal numbers]
  \label{def:surreal-number-mul}
  \lean{Surreal.SurrealNumber.mul,Surreal.instMulSurrealNumber}
  \leanok
  \uses{def:surreal-number,thm:surreal-mul}
  Define multiplication on \(\NoShort\) by taking the game product of any two representatives
  and then passing to the quotient.
\end{definition}

\begin{theorem}[Multiplicative laws on the quotient]
  \label{thm:quotient-mul-laws}
  \lean{Surreal.SurrealNumber.mul_comm,Surreal.SurrealNumber.mul_assoc,Surreal.SurrealNumber.one_mul,Surreal.SurrealNumber.mul_one,Surreal.SurrealNumber.zero_mul,Surreal.SurrealNumber.mul_zero}
  \leanok
  Multiplication on \(\NoShort\) is associative and commutative, has identity \(1\), and has
  \(0\) as an absorbing element.
\end{theorem}
\begin{proof}
  \leanok
  \uses{def:surreal-number,def:surreal-number-mul,thm:game-mul-laws,thm:game-mul-assoc}
  By Definition~\ref{def:surreal-number}, choose surreal representatives
  \(a_0,b_0,c_0\) for three arbitrary classes.  Definition~\ref{def:surreal-number-mul}
  identifies their products with the classes of the corresponding game
  products.  The commutativity, unit, and zero equivalences of
  Theorem~\ref{thm:game-mul-laws} therefore give
  \[
    [a_0b_0]=[b_0a_0],\qquad [1\cdot a_0]=[a_0]=[a_0\cdot1],
    \qquad [0\cdot a_0]=[0]=[a_0\cdot0].
  \]
  Likewise, Theorem~\ref{thm:game-mul-assoc}, applied to the surreal games
  \(a_0,b_0,c_0\), gives
  \[
    [(a_0b_0)c_0]=[a_0(b_0c_0)].
  \]
  Equality of each displayed pair of classes follows precisely because the
  representative games are equivalent.  Since the representatives were
  arbitrary, all the asserted laws hold on \(\NoShort\).
\end{proof}

\begin{theorem}[Distributivity on the quotient]
  \label{thm:quotient-distrib}
  \lean{Surreal.SurrealNumber.left_distrib,Surreal.SurrealNumber.right_distrib}
  \leanok
  For \(a,b,c\in\NoShort\),
  \[
    a(b+c)=ab+ac,
    \qquad
    (a+b)c=ac+bc.
  \]
\end{theorem}
\begin{proof}
  \leanok
  \uses{def:surreal-number,thm:game-distrib,def:surreal-number-mul,thm:quotient-mul-laws}
  Choose surreal representatives \(a_0,b_0,c_0\) as permitted by
  Definition~\ref{def:surreal-number}.  By
  Definition~\ref{def:surreal-number-mul}, left distributivity reduces to the
  equivalence
  \[
    a_0(b_0+c_0)\GameEq a_0b_0+a_0c_0,
  \]
  which is Theorem~\ref{thm:game-distrib}.  Equivalent representatives define
  the same class, so this proves \(a(b+c)=ab+ac\).

  For right distributivity, use the commutativity in
  Theorem~\ref{thm:quotient-mul-laws} to write
  \[
    (a+b)c=c(a+b)=ca+cb=ac+bc,
  \]
  where the middle equality is the left distributive law just proved and the
  last equality again uses commutativity.
\end{proof}

\begin{theorem}[Commutative ring structure]
  \label{thm:comm-ring}
  \lean{Surreal.instCommRingSurrealNumber}
  \leanok
  The short surreal numbers \(\NoShort\) form a commutative ring.
\end{theorem}
\begin{proof}
  \leanok
  \uses{thm:ordered-add-group,def:surreal-number-mul,thm:quotient-mul-laws,thm:quotient-distrib}
  Theorem~\ref{thm:ordered-add-group} supplies the additive commutative group,
  including associativity, commutativity, the zero element, and additive
  inverses.  Definition~\ref{def:surreal-number-mul} supplies multiplication.
  Its associativity, commutativity, identity element, and zero laws are exactly
  Theorem~\ref{thm:quotient-mul-laws}, while both distributive axioms are
  Theorem~\ref{thm:quotient-distrib}.  These are precisely the axioms of a
  commutative ring.
\end{proof}

\section{Positive products and ordered-ring compatibility}

\begin{theorem}[A product of positive surreal games is positive]
  \label{thm:game-mul-pos}
  \lean{Surreal.Game.mul_pos_of_isSurreal}
  \leanok
  If \(x,y\) are surreal games and \(0\GameLt x\), \(0\GameLt y\), then
  \[
    0\GameLt xy.
  \]
\end{theorem}
\begin{proof}
  \leanok
  \uses{def:game,def:game-order,def:game-inductions,lem:birthday-options,def:is-surreal,def:cross-conditions,thm:game-order-calculus,thm:game-add-laws,thm:game-mul-laws,thm:conway-b,thm:conway-c}
  First show that \(0\GameLt y\) forces a left option \(y^L\) with
  \(0\GameLe y^L\).  If no left option had this property, the left clause in
  Definition~\ref{def:game-order} would give \(y\GameLe0\); the other clause
  is vacuous because \(0\) has no right options.  This contradicts the strict
  hypothesis.

  Proceed by the well-founded scheme of
  Definition~\ref{def:game-inductions}, using
  Lemma~\ref{lem:birthday-options}, and choose such a \(y^L\).  By
  Definition~\ref{def:is-surreal}, the option \(y^L\) is surreal.  Split according as
  \(y^L\GameLe0\).  If it does, then the two weak comparisons give
  \(y^L\GameEq0\).  Theorem~\ref{thm:conway-b}, together with the
  commutativity and zero-product laws in Theorem~\ref{thm:game-mul-laws}, gives
  \(xy^L\GameEq0\), hence \(0\GameLe xy^L\).  Otherwise
  Definition~\ref{def:game-order}, together with the already known comparison
  \(0\GameLe y^L\), gives \(0\GameLt y^L\).  Since \(y^L\) has smaller birthday, the induction
  hypothesis gives \(0\GameLt xy^L\), and therefore again
  \(0\GameLe xy^L\).

  Apply Theorem~\ref{thm:conway-c} to \(0\GameLt x\), with common surreal
  factor \(y\) and left option \(y^L\).  Simplifying its left cross inequality
  by the zero-product laws of Theorem~\ref{thm:game-mul-laws} and the zero-sum
  laws of Theorem~\ref{thm:game-add-laws} yields
  \(xy^L\GameLt xy\).  The mixed transitivity assertion of
  Theorem~\ref{thm:game-order-calculus}, applied to
  \(0\GameLe xy^L\GameLt xy\), proves \(0\GameLt xy\).
\end{proof}

\begin{theorem}[Positive and nonnegative products on the quotient]
  \label{thm:quotient-mul-pos}
  \lean{Surreal.SurrealNumber.zero_le_one,Surreal.SurrealNumber.mul_pos,Surreal.SurrealNumber.mul_nonneg,Surreal.SurrealNumber.instZeroLEOneClass}
  \leanok
  For \(a,b\in\NoShort\),
  \[
    0<a,\ 0<b\Longrightarrow 0<ab,
    \qquad
    0\leq a,\ 0\leq b\Longrightarrow 0\leq ab,
  \]
  and \(0\leq1\).
\end{theorem}
\begin{proof}
  \leanok
  \uses{def:game,def:game-order,def:surreal-number,thm:ordered-add-group,def:surreal-number-mul,thm:quotient-mul-laws,thm:game-mul-pos}
  For strict positivity, choose surreal representatives \(a_0,b_0\) of the
  two classes using Definition~\ref{def:surreal-number}.  The hypotheses become
  \(0\GameLt a_0\) and \(0\GameLt b_0\).  By
  Theorem~\ref{thm:game-mul-pos}, \(0\GameLt a_0b_0\); passing to equivalence
  classes and using Definition~\ref{def:surreal-number-mul} gives \(0<ab\).

  Suppose only \(0\leq a\) and \(0\leq b\).  The linear order of
  Theorem~\ref{thm:ordered-add-group} splits each weak comparison into a strict
  inequality or equality.  If both inequalities are strict, the first part of
  this theorem gives \(0<ab\), hence \(0\leq ab\).  If \(a=0\) or \(b=0\),
  the appropriate zero-product law in Theorem~\ref{thm:quotient-mul-laws}
  gives \(ab=0\).  These cases prove nonnegative-product closure.

  Finally, Definitions~\ref{def:game} and~\ref{def:game-order} show directly
  that \(0\GameLe1\): the left-option condition for \(0\) and the right-option
  condition for \(1\) are both vacuous.  Hence their equivalence classes
  satisfy \(0\leq1\).
\end{proof}

\begin{theorem}[Monotonicity of multiplication]
  \label{thm:quotient-mul-mono}
  \lean{Surreal.SurrealNumber.mul_le_mul_of_nonneg_left',Surreal.SurrealNumber.mul_le_mul_of_nonneg_right',Surreal.SurrealNumber.mul_lt_mul_of_pos_left',Surreal.SurrealNumber.mul_lt_mul_of_pos_right',Surreal.SurrealNumber.instPosMulMono,Surreal.SurrealNumber.instMulPosMono,Surreal.SurrealNumber.instPosMulStrictMono,Surreal.SurrealNumber.instMulPosStrictMono}
  \leanok
  Multiplication by a nonnegative short surreal number preserves weak inequalities on either
  side, and multiplication by a positive short surreal number preserves strict inequalities on
  either side.
\end{theorem}
\begin{proof}
  \leanok
  \uses{thm:ordered-add-group,thm:comm-ring,thm:quotient-mul-pos}
  Assume \(b\leq c\) and \(0\leq a\).  By
  Theorem~\ref{thm:ordered-add-group}, the first hypothesis is equivalent to
  \(0\leq c-b\).  The nonnegative-product assertion of
  Theorem~\ref{thm:quotient-mul-pos} gives \(0\leq a(c-b)\), while the ring
  laws of Theorem~\ref{thm:comm-ring} give \(a(c-b)=ac-ab\).  The ordered
  additive group equivalence
  \(0\leq ac-ab\Longleftrightarrow ab\leq ac\) yields weak monotonicity on
  the left.  For the right-sided assertion, suppose \(a\leq b\) and
  \(0\leq c\).  Then \(0\leq b-a\), so
  \[
    0\leq(b-a)c=bc-ac,
  \]
  and hence \(ac\leq bc\).

  If \(b<c\) and \(0<a\), use instead \(0<c-b\).  Strict positivity of a
  product, again by Theorem~\ref{thm:quotient-mul-pos}, gives
  \(0<a(c-b)=ac-ab\), which Theorem~\ref{thm:ordered-add-group} identifies
  with \(ab<ac\).  Similarly, \(a<b\) and \(0<c\) imply
  \[
    0<(b-a)c=bc-ac,
  \]
  and hence \(ac<bc\).  These four implications give the four weak and strict,
  left- and right-sided multiplication-order compatibilities asserted in the
  theorem.
\end{proof}

\begin{theorem}[Ordered- and strictly ordered-ring compatibility]
  \label{thm:ordered-ring}
  \lean{Surreal.instIsOrderedRingSurrealNumber,Surreal.instIsStrictOrderedRingSurrealNumber}
  \leanok
  The ring \(\NoShort\) is nontrivial; addition reflects order, multiplication by a
  nonnegative element preserves weak inequalities, and multiplication by a positive element
  preserves strict inequalities, on either side.
\end{theorem}
\begin{proof}
  \leanok
  \uses{def:game,def:game-order,def:surreal-number,thm:game-order-calculus,thm:ordered-add-group,thm:comm-ring,thm:quotient-mul-mono}
  Theorem~\ref{thm:ordered-add-group} gives translation invariance of the
  order, Theorem~\ref{thm:comm-ring} gives the commutative ring laws, and the
  weak and strict multiplication compatibility assertions are
  Theorem~\ref{thm:quotient-mul-mono}.  In particular, if
  \(c+a\leq c+b\), translating both sides by \(-c\) and simplifying with the
  additive group laws of Theorem~\ref{thm:ordered-add-group} gives
  \(a\leq b\).  Thus addition reflects order as required.

  It remains to supply nontriviality.  At game level, \(0\) is the sole left
  option of \(1\) by Definition~\ref{def:game}.  If \(1\GameLe0\), the left
  clause of Definition~\ref{def:game-order} would require
  \(\neg(0\GameLe0)\), contradicting the reflexivity assertion of
  Theorem~\ref{thm:game-order-calculus}.  Therefore \(0\) and \(1\) are not
  equivalent, so their classes in Definition~\ref{def:surreal-number} are
  distinct.  Order reflection and the weak and strict multiplication
  inequalities above establish all the claimed ordered-ring
  compatibility axioms.
\end{proof}

\begin{corollary}[Main endpoint: a linear ordered commutative ring]
  \label{thm:main-endpoint}
  \lean{Surreal.instLinearOrderSurrealNumber,Surreal.instCommRingSurrealNumber,Surreal.instIsStrictOrderedRingSurrealNumber}
  \leanok
  The short surreal numbers constructed above form a linear ordered commutative ring.
\end{corollary}
\begin{proof}
  \leanok
  \uses{thm:ordered-add-group,thm:comm-ring,thm:ordered-ring}
  By Theorem~\ref{thm:ordered-add-group}, \(\NoShort\) is a linear ordered
  additive commutative group.  By Theorem~\ref{thm:comm-ring}, \(\NoShort\) is
  a commutative ring.
  Theorem~\ref{thm:ordered-ring} proves compatibility of these algebraic
  operations with the linear order.  Hence \(\NoShort\) is a linear ordered
  commutative ring.
\end{proof}
